Для функции
\[
f(x_1,x_2)=f(X), \qquad X=(x_1,x_2),
\]
на \(k\)-м шаге задаётся текущее приближение
\[
X^k=(x_1^k,x_2^k).
\]

Далее рассматривается функция одной переменной
\[
\varphi_k(t)=f\bigl(X^k-t\,\mathrm{grad}\,f(X^k)\bigr).
\]
Здесь градиент функции в точке \(X^k\):
\[
\mathrm{grad}\,f(X^k)=
\left(
\frac{\partial f}{\partial x_1}(X^k),
\frac{\partial f}{\partial x_2}(X^k)
\right)
\]

Функция \(\varphi_k(t)\) показывает, как изменяется \(f\) при движении из точки \(X^k\) в направлении, противоположном градиенту, то есть в направлении наискорейшего убывания.

Следующее приближение определяется по формуле
\[
X^{k+1}=X^k-t^*\,\mathrm{grad}\,f(X^k),
\]
или в координатной форме
\[
x_1^{k+1}=x_1^k-t^*\frac{\partial f}{\partial x_1}(X^k),
\qquad
x_2^{k+1}=x_2^k-t^*\frac{\partial f}{\partial x_2}(X^k).
\]

\subsection*{Шаг спуска}

Так как функция зависит от двух переменных, а движение на каждом шаге идёт вдоль фиксированного направления \(-\mathrm{grad}\,f(X^k)\), задача выбора шага сводится к минимизации функции одной переменной \(\varphi_k(t)\).

Шаг спуска находится приближённо по методу парабол:
\[
t^*=-\frac{\varphi_k'(0)}{\varphi_k''(0)}.
\]

Для нашего двумерного случая:
\[
(x_{k+1},y_{k+1})=
\left(
x_k-t^*\frac{\partial f}{\partial x},
\;
y_k-t^*\frac{\partial f}{\partial y}
\right),
\]
где
\[
\varphi_k'(0)=-
\left(\frac{\partial f}{\partial x}\right)^2-
\left(\frac{\partial f}{\partial y}\right)^2,
\]
\[
\varphi_k''(0)=
\frac{\partial^2 f}{\partial x^2}\left(\frac{\partial f}{\partial x}\right)^2
+
2\frac{\partial^2 f}{\partial x\partial y}
\frac{\partial f}{\partial x}
\frac{\partial f}{\partial y}
+
\frac{\partial^2 f}{\partial y^2}\left(\frac{\partial f}{\partial y}\right)^2.
\]

Все производные в этих формулах вычисляются в текущей точке \((x_k,y_k)\).

\subsection*{Критерий остановки}

Итерации продолжаются, пока не выполнится условие
\[
\|\mathrm{grad}\,f(X^k)\|<\varepsilon.
\]

Здесь норма градиента определяется как максимум по абсолютным значениям его компонент:
\[
\|\mathrm{grad}\,f(X^k)\|=
\max\left(
\left|\frac{\partial f}{\partial x_1}(X^k)\right|,
\left|\frac{\partial f}{\partial x_2}(X^k)\right|
\right).
\]

\subsection*{Аналитическое нахождение минимума}

Рассмотрим нашу функцию
\[
f(x_1,x_2)=2x_1^2+3x_2^2-2\sin\left(\frac{x_1-x_2}{2}\right)+x_2.
\]

Чтобы найти стационарную точку, приравниваем к нулю частные производные:
\[
\frac{\partial f}{\partial x_1}=4x_1-\cos\left(\frac{x_1-x_2}{2}\right)=0,
\]
\[
\frac{\partial f}{\partial x_2}=6x_2+\cos\left(\frac{x_1-x_2}{2}\right)+1=0.
\]

Отсюда получаем
\[
x_1=\frac{1}{4}\cos\left(\frac{x_1-x_2}{2}\right),
\qquad
x_2=-\frac{1+\cos\left(\frac{x_1-x_2}{2}\right)}{6}.
\]

Для упрощения вводим обозначение

\[
u=\frac{x_1-x_2}{2}.
\]

Тогда
\[
x_1=\frac{\cos u}{4},
\qquad
x_2=-\frac{1+\cos u}{6}.
\]

Подставим эти выражения в формулу для \(u\):
\[
u=\frac{x_1-x_2}{2}
=
\frac{1}{2}
\left(
\frac{\cos u}{4}
+
\frac{1+\cos u}{6}
\right).
\]

Приводя к общему знаменателю, получаем уравнение
\[
u=\frac{5\cos u+2}{24}.
\]

Это уравнение решается численно. Получаем
\[
u \approx 0.2833587125.
\]

Тогда координаты точки минимума:
\[
x_1^*=\frac{\cos u}{4}\approx 0.240030,
\qquad
x_2^*=-\frac{1+\cos u}{6}\approx -0.326687.
\]

Минимальное значение функции:
\[
f_{\min}=f(x_1^*,x_2^*)\approx -0.450449.
\]


Таким образом, аналитическое решение получается из системы уравнений
\[
\frac{\partial f}{\partial x_1}=0,
\qquad
\frac{\partial f}{\partial x_2}=0,
\]
после чего вводится замена
\[
u=\frac{x_1-x_2}{2},
\]
и через найденное значение \(u\) вычисляются координаты точки минимума и значение функции в этой точке.

